\chapter{Cadre théorique} \label{chap:cadre_theoriques}

Introduction et annonce du plan du chapitre.

\lipsum[1]

\section{Section 1} \label{sec:titre_cth_section1}

Paragraphe.
\lipsum[1] % génère un paragraphe (à supprimer)


\section{Section 2} \label{sec:titre_cth_section2}

Paragraphe.
\lipsum[1] % génère un paragraphe (à supprimer)

\textbf{Comment ajouter une courte citation?} Selon \textcite[69]{bonn1992precis}, «~le comportement spectral de la végétation diffère sensiblement de celui des sols et des roches~».

\textbf{Comment ajouter une citation de plus de trois lignes?}

\begin{customquote}
« Le comportement spectral de la végétation diffère sensiblement de celui des sols et des roches. La végétation est un milieu complexe et changeant dans le temps, dont les propriétés spectrales varient avec la saison et les phases de croissance. Les études de propriétés spectrales de la végétation ont été réalisées à différentes échelles, dont les principales sont celle de la feuille, celle de la plante et celle du couvert, par ordre de complexité croissante » \parencite[69]{bonn1992precis}. 
\end{customquote}

Plusieurs études menées au Canada ont démontré que... \parencite{manaugh2017overcoming,winters2011motivators}.

\textbf{Connection avec Zotero}

Si vous utilisez \href{https://www.overleaf.com}{Overleaf}, il est possible de le connecter à votre compte 
\href{https://docs.overleaf.com/integrations-and-add-ons/reference-manager-integrations/zotero}{Zotero}.

\section{Section 3} \label{sec:titre_cth_section3}

\lipsum[1-4] % génère quatre paragraphes (à supprimer)
